\chapter{Osteoporosis}

\section*{Definition and Overview}
Osteoporosis is a systemic skeletal disease characterised by low bone mass and microarchitectural deterioration of bone tissue, leading to increased bone fragility and fracture risk. It is defined operationally by bone mineral density (BMD) T-score $\le$--2.5 standard deviations below the young adult mean at the femoral neck or lumbar spine (DXA). Osteoporotic fractures (minimal trauma fractures) occur most commonly at the hip, spine, wrist, and humerus. Osteoporosis is often asymptomatic until a fracture occurs, earning it the name ``silent disease.'' It is highly prevalent, treatable, and largely preventable.

\section*{Epidemiology}
In Australia, an estimated 1.2 million people have osteoporosis (T-score $\le$--2.5) and 6.3 million have osteopenia (T-score --1.0 to --2.5). One in three women and one in five men over 50 will sustain an osteoporotic fracture. Hip fracture incidence is ~20,000/year; vertebral fractures are more common but often undiagnosed. Fracture risk increases exponentially with age; women lose bone rapidly after menopause (10-year accelerated loss). Indigenous Australians have lower BMD but higher fracture rates at younger ages. Annual cost exceeds \$3.8 billion. Despite effective treatments, a ``treatment gap'' persists: <30\% of high-risk patients receive therapy.

\section*{Aetiology and Pathophysiology}
\begin{itemize}
    \item \textbf{Bone remodelling:} continuous cycle of osteoclast-mediated resorption and osteoblast-mediated formation; coupled in health, uncoupled in osteoporosis (resorption > formation)
    \item \textbf{Peak bone mass:} achieved by age 20--30; determined by genetics (60--80\%), nutrition, physical activity, hormonal status; low peak mass increases lifetime risk
    \item \textbf{Menopause:} oestrogen deficiency $\to$ increased osteoclast activity, decreased osteoblast lifespan, RANKL/OPG imbalance; rapid trabecular bone loss (spine) then cortical (hip)
    \item \textbf{Ageing:} reduced osteoblast function, increased marrow adiposity, oxidative stress, senescent cell accumulation; cortical porosity, endocortical resorption
    \item \textbf{Secondary causes (always exclude):}
    \begin{itemize}
        \item Endocrine: hyperthyroidism, hyperparathyroidism, Cushing, hypogonadism (men), type 1 diabetes, premature menopause
        \item GI: coeliac disease, IBD, malabsorption, bariatric surgery
        \item Renal: CKD-MBD
        \item Haematological: myeloma, lymphoma
        \item Medications: glucocorticoids ($\ge$5 mg prednisolone $\ge$3 months), aromatase inhibitors, androgen deprivation therapy, PPIs (long-term), anticonvulsants, heparin, SSRIs
        \item Lifestyle: smoking, alcohol excess, low calcium/vitamin D, inactivity, low body weight (BMI <20)
    \end{itemize}
    \item \textbf{Genetics:} multiple genes (LRP5, WNT16, RANKL, OPG, COL1A1, VDR); polygenic risk scores emerging
\end{itemize}

\section*{Clinical Features}
\begin{itemize}
    \item \textbf{Asymptomatic:} most patients until fracture
    \item \textbf{Fractures:}
    \begin{itemize}
        \item Vertebral: often silent (2/3); height loss $>$3 cm, kyphosis, back pain
        \item Hip: fall from standing height; inability to bear weight; shortened/externally rotated leg
        \item Distal radius (Colles'): fall on outstretched hand
        \item Proximal humerus, pelvis, ribs
    \end{itemize}
    \item \textbf{Signs:} height loss, thoracic kyphosis (``dowager's hump''), reduced occiput-to-wall distance, tenderness over spinous processes
    \item \textbf{Red flags for secondary cause:} early onset (<50 men, premenopausal women), severe osteoporosis (Z-score $\le$--2.0), atypical fractures (femoral shaft, stress fractures), multiple fractures
\end{itemize}

\section*{Differential Diagnosis}
\begin{itemize}
    \item \textbf{Osteomalacia:} defective mineralisation; low calcium/phosphate, high ALP, high PTH, low 25-OH vitamin D; Looser zones (pseudofractures)
    \item \textbf{Paget disease:} focal increased turnover; high ALP, sclerotic/lucent bone on imaging; usually asymptomatic
    \item \textbf{Multiple myeloma:} lytic lesions, anaemia, renal impairment, high ESR, M-protein
    \item \textbf{Metastatic bone disease:} known primary, multiple lesions, often blastic (prostate, breast) or lytic (lung, kidney, thyroid)
    \item \textbf{Osteogenesis imperfecta:} childhood fractures, blue sclerae, dentinogenesis imperfecta, hearing loss
    \item \textbf{Idiopathic juvenile osteoporosis:} rare, self-limiting
    \item \textbf{Medication effects:} glucocorticoids, aromatase inhibitors, ADT
\end{itemize}

\section*{When to Seek Medical Care}
Assess fracture risk in all adults from age 50 (women) or 60 (men), or earlier with risk factors. Consider DXA if: women $\ge$65, men $\ge$70, postmenopausal women <65 with risk factors, men 50--69 with risk factors, adults with fragility fracture, glucocorticoid therapy $\ge$3 months, high fall risk.

\textbf{Urgent:}
\begin{itemize}
    \item Suspected hip fracture (fall + hip pain + inability to weight-bear)
    \item Vertebral fracture with neurological symptoms (cord compression)
    \item Severe back pain unresponsive to analgesia
\end{itemize}

\section*{Diagnosis and Investigations}
\begin{itemize}
    \item \textbf{DEXA (DXA):} gold standard; femoral neck and lumbar spine (L1--L4); T-score for postmenopausal women and men $\ge$50; Z-score for premenopausal women, men <50, children; lateral vertebral imaging (VFA) for silent vertebral fractures
    \item \textbf{Fracture risk calculators:} FRAX (10-year major osteoporotic and hip fracture probability); Garvan Fracture Risk Calculator (includes falls history); Australian FRAX model available
    \item \textbf{Bloods (exclude secondary causes):} calcium, phosphate, ALP, 25-OH vitamin D, PTH, TSH, creatinine, eGFR, FBC, ESR, coeliac serology (tTG-IgA), testosterone (men), 24-h urinary calcium
    \item \textbf{Bone turnover markers (optional):} CTX (resorption), P1NP (formation); monitor treatment response
    \item \textbf{Imaging:} X-ray thoracic/lumbar spine for vertebral fracture assessment (Genant semiquantitative); MRI if acute vs chronic fracture uncertain; CT for complex fractures
    \item \textbf{Falls assessment:} gait, balance, vision, medications, home hazards, cardiovascular, neurological
\end{itemize}

\section*{Management}
\begin{itemize}
    \item \textbf{Non-pharmacological (all patients):}
    \begin{itemize}
        \item Calcium: 1000 mg/d (19--50 yr), 1300 mg/d (women >50, men >70); dietary preferred (dairy, fortified, leafy greens, fish with bones); supplement if dietary insufficient
        \item Vitamin D: target 25-OH D $\ge$50 nmol/L (end winter); 1000--2000 IU/d typical; higher doses if deficient
        \item Exercise: weight-bearing aerobic (walking, jogging), progressive resistance training (2--3x/week), balance training (Tai Chi); avoid high-impact if severe osteoporosis
        \item Falls prevention: home safety assessment, medication review (sedatives, antihypertensives), vision correction, footwear, hip protectors for high-risk
        \item Smoking cessation, alcohol moderation ($\le$2 standard drinks/day)
        \item Protein: $\ge$1.0--1.2 g/kg/day
    \end{itemize}
    \item \textbf{Pharmacological (T-score $\le$--2.5, or fragility fracture, or high FRAX risk):}
    \begin{itemize}
        \item \textbf{First-line:} Oral bisphosphonates (alendronate 70 mg weekly, risedronate 35 mg weekly); IV zoledronic acid 5 mg annually (if oral intolerant/contraindicated); PBS criteria: T-score $\le$--2.5 + age $\ge$70, or fragility fracture, or glucocorticoid use
        \item \textbf{Denosumab:} 60 mg SC 6-monthly; potent antiresorptive; must not delay doses (rebound vertebral fractures); monitor calcium/vitamin D; hypocalcaemia risk in CKD
        \item \textbf{Anabolic agents (high fracture risk, severe osteoporosis, failed antiresorptives):} Teriparatide 20 $\mu$g SC daily $\times$ 24 months (PBS: $\ge$2 fragility fractures, one hip/spine, T-score $\le$--3.0); Romosozumab 210 mg SC monthly $\times$ 12 months (PBS: similar criteria) followed by antiresorptive
        \item \textbf{HRT:} oestrogen +/- progestogen for menopausal women <60 or <10 yr post-menopause with indications; prevents fractures but not first-line for osteoporosis alone
        \item \textbf{Raloxifene:} SERM; reduces vertebral fractures; increases VTE risk; consider for postmenopausal women with breast cancer risk
        \item \textbf{Glucocorticoid-induced:} bisphosphonate or denosumab if $\ge$5 mg pred $\ge$3 months; calcium/vitamin D; consider teriparatide for high risk
    \end{itemize}
    \item \textbf{Monitoring:} DEXA at 1--2 years after starting/changing therapy; bone turnover markers at 3--6 months; annual height, falls review; adherence counselling
    \item \textbf{Treatment holiday:} after 3--5 years oral bisphosphonate or 3 years zoledronic acid if low fracture risk (T-score >--2.5, no fractures); monitor annually; restart if fracture or significant BMD loss; denosumab -- NO holiday (transition to bisphosphonate)
\end{itemize}

\section*{Complications}
\begin{itemize}
    \item \textbf{Hip fracture:} 20--30\% 1-year mortality; 50\% lose independent mobility; 20\% long-term care
    \item \textbf{Vertebral fracture:} chronic pain, height loss, kyphosis, restrictive lung disease, reduced QoL; 5x risk subsequent vertebral fracture
    \item \textbf{Refracture:} 30--50\% within 5 years; highest risk first 1--2 years post-fracture (``imminent risk'')
    \item \textbf{Medication adverse effects:} bisphosphonates -- oesophagitis, atypical femoral fracture (rare, long-term), ONJ (very rare, dental); denosumab -- hypocalcaemia, rebound fractures, ONJ; anabolics -- hypercalcaemia, osteosarcoma risk (theoretical, rat data)
    \item \textbf{Psychosocial:} fear of falling, activity restriction, depression, loss of independence
\end{itemize}

\section*{Prognosis and Outcomes}
With treatment, vertebral fracture risk reduced 40--70\%, hip fracture risk 40--50\%. Anabolics build bone (spine BMD +10--15\% in 12--18 months). Fracture liaison services (FLS) improve identification and treatment initiation post-fracture by 3--4x. Early diagnosis and consistent therapy preserve independence and reduce mortality. The treatment gap remains the major challenge.

\section*{Prevention and Self-Care}
\begin{itemize}
    \item Build peak bone mass in youth: calcium, vitamin D, weight-bearing exercise, avoid smoking/alcohol
    \item Maintain adequate calcium (diet first) and vitamin D (sun exposure: 5--10 min mid-morning/afternoon arms exposed; supplement if needed)
    \item Lifelong weight-bearing and resistance exercise
    \item Falls-proof home: lighting, grab rails, non-slip mats, remove clutter
    \item Regular medication review (deprescribe sedatives, anticholinergics)
    \item Bone health check at menopause (women) and age 70 (men), or earlier with risk factors
    \item If on treatment: take as directed (bisphosphonates: fasting, upright 30 min); don't stop without doctor advice; denosumab -- never miss a dose
    \item Calcium/vitamin D supplements with every osteoporosis medication
    \item Join support: Healthy Bones Australia, Osteoporosis Australia
\end{itemize}

\section*{Sources and Further Reading}
The following reputable sources provide further information for healthcare students and patients seeking reliable, current advice about osteoporosis.

\begin{itemize}
    \item Healthy Bones Australia (formerly Osteoporosis Australia). \textit{Clinical guidelines and patient resources.} https://www.healthybonesaustralia.org.au/
    \item RACGP. \textit{Osteoporosis prevention, diagnosis and management.} https://www.racgp.org.au/
    \item Royal Australian College of Physicians. \textit{Osteoporosis clinical guidelines.} https://www.racp.edu.au/
    \item Endocrine Society of Australia. \textit{Position statements on osteoporosis.} https://www.endocrinesociety.org.au/
    \item Australian and New Zealand Bone and Mineral Society (ANZBMS). \textit{Fracture liaison services and guidelines.} https://www.anzbms.org.au/
    \item Healthdirect Australia. \textit{Osteoporosis.} https://www.healthdirect.gov.au/osteoporosis
    \item International Osteoporosis Foundation (IOF). \textit{Global guidelines and fracture risk calculators.} https://www.osteoporosis.foundation/
    \item National Institute for Health and Care Excellence (NICE). \textit{Osteoporosis guidelines.} https://www.nice.org.uk/
\end{itemize}